<<<<<<< HEAD
\subsection{BÀI TẬP TRẮC NGHIỆM}

\ind{PHẦN I.} \inden{Câu trắc nghiệm nhiều phương án lựa chọn. Học sinh trả lời từ câu 1 đến câu 20. Mỗi câu hỏi học sinh chỉ chọn một phương án.}\\
\setcounter{ex}{0}
\Opensolutionfile{ans}[ans/1D1-B2-1]
\begin{ex}%[0D6Y3-8]
	Khẳng định nào sau đây là khẳng định \textbf{sai}?
	\choice
	{$\cos \left({a-b}\right)=\sin a\sin b+\cos a\cos b$}
	{\True $\cos \left({a+b}\right)=\sin a\sin b-\cos a\cos b$}
	{$\sin \left({a-b}\right)=\sin a\cos b-\cos a\sin b$}
	{$\sin \left({a+b}\right)=\sin a\cos b+\cos a\sin b$}
	\loigiai{Ta có $\cos \left({a+b}\right)=\cos a\cos b-\sin a\sin b$. } 
\end{ex}\cham{2}

\begin{ex}%[0D6B3-8]
	Khẳng định nào sau đây là khẳng định đúng?
	\choice
	{$\sin a+\cos a=\sqrt{2}\sin \left({a-\dfrac{\pi}{4}}\right)$}
	{\True $\sin a+\cos a=\sqrt{2}\sin \left({a+\dfrac{\pi}{4}}\right)$}
	{$\sin a+\cos a=-\sqrt{2}\sin \left({a-\dfrac{\pi}{4}}\right)$}
	{$\sin a+\cos a=-\sqrt{2}\sin \left({a+\dfrac{\pi}{4}}\right)$}
	\loigiai{  } 
\end{ex}\cham{2}

\begin{ex}%[0D6B3-1]%
	Cho các góc lượng giác $a,b$ và $T=\cos (a+b)\cos (a-b)-\sin (a+b)\sin (a-b)$. Mệnh đề sau đây đúng?
	\choice
	{$T=\sin 2b$}
	{$T=\sin 2a$}
	{\True $T=\cos 2a$}
	{$T=\cos 2b$}
	\loigiai{
		Ta có $T=\cos (a+b)\cos (a-b)-\sin (a+b)\sin (a-b)=\cos [(a+b)+(a-b)]=\cos 2a$.}
\end{ex}\cham{6}

\begin{ex}%[0D6B3-1]%
	Cho $\sin \alpha=\dfrac{1}{\sqrt{3}}$ với $0<\alpha<\dfrac{\pi}{2}$. Khi đó, giá trị của $\sin \left(\alpha+\dfrac{\pi}{3}\right)$ bằng
	\choice
	{$\dfrac{\sqrt{3}}{3}-\dfrac{1}{2}$}
	{\True $\dfrac{\sqrt{3}}{6}+\dfrac{\sqrt{2}}{2}$}
	{$\sqrt{6}-\dfrac{1}{2}$}
	{$\dfrac{\sqrt{3}}{6}-\dfrac{\sqrt{2}}{2}$}
	\loigiai{
		Với $0<\alpha<\dfrac{\pi}{2}$ thì $\cos\alpha >0$ nên $\cos \alpha=\sqrt{1-\sin^2 \alpha}=\dfrac{\sqrt{6}}{3}$.\\
		Ta có $\sin \left(\alpha+\dfrac{\pi}{3}\right)=\sin \alpha\cdot \cos \dfrac{\pi}{3}+\cos \alpha\cdot \sin \dfrac{\pi}{3}=\dfrac{1}{\sqrt{3}}\cdot \dfrac{1}{2}+\dfrac{\sqrt{6}}{3}\cdot \dfrac{\sqrt{3}}{2}=\dfrac{\sqrt{3}}{6}+\dfrac{\sqrt{2}}{2}$.
	}
\end{ex}\cham{6}


\begin{ex}%[0D6B3-4] 
	Biết $\sin \alpha=-\dfrac{3}{5}$ và $\pi <\alpha <\dfrac{{3\pi}}{2}$. Tính $\sin \left({\alpha+\dfrac{\pi}{6}}\right).$
	\choice
	{$-\dfrac{3}{5}$}
	{$\dfrac{3}{5}$}
	{\True $\dfrac{{-4-3\sqrt{3}}}{{10}}$}
	{$\dfrac{{4-3\sqrt{3}}}{{10}}$}
	\loigiai{
		Từ hệ thức $\sin ^2\alpha+\cos ^2\alpha =1$, suy ra $\cos \alpha =\pm \sqrt{{1-{\sin}^2\alpha}}=\pm \dfrac{4}{5}$.
		Do $\pi <\alpha <\dfrac{{3\pi}}{2}$ nên ta chọn $\cos \alpha =-\dfrac{4}{5}$.\\
		Suy ra $P=\sin \left({\alpha+\dfrac{\pi}{6}}\right)=\dfrac{{\sqrt{3}}}{2}\sin \alpha+\dfrac{1}{2}\cos \alpha =\dfrac{{\sqrt{3}}}{2}\left({-\dfrac{3}{5}}\right)+\dfrac{1}{2}\left({-\dfrac{4}{5}}\right)=\dfrac{{-4-3\sqrt{3}}}{{10}}$.
	} 
\end{ex}\cham{7}

\begin{ex}%[0D6B3-4] 
	Cho góc $\alpha $ thỏa mãn $\cos \alpha =\dfrac{3}{4}$ và $\dfrac{{3\pi}}{2}<\alpha <2\pi $. Tính $\cos \left({\dfrac{\pi}{3}-\alpha}\right).$
	\choice
	{$\dfrac{{3+\sqrt{{21}}}}{8}$}
	{\True $\dfrac{{3-\sqrt{{21}}}}{8}$}
	{$\dfrac{{3\sqrt{3}+\sqrt{7}}}{8}$}
	{$\dfrac{{3\sqrt{3}-\sqrt{7}}}{8}$}
	\loigiai{ Ta có $P=\cos \left({\dfrac{\pi}{3}-\alpha}\right)=\cos \dfrac{\pi}{3}\cos \alpha+\sin \dfrac{\pi}{3}\sin \alpha =\dfrac{1}{2}\cos \alpha+\dfrac{{\sqrt{3}}}{2}\sin \alpha $.\\
		Từ hệ thức $\sin ^2\alpha+\cos ^2\alpha =1$, suy ra $\sin \alpha =\pm \sqrt{{1-{\cos}^2\alpha}}=\pm \dfrac{{\sqrt{7}}}{4}$.\\
		Do $\dfrac{{3\pi}}{2}<\alpha <2\pi $ nên ta chọn $\sin \alpha =-\dfrac{{\sqrt{7}}}{4}$.\\
		Thay $\sin \alpha =-\dfrac{{\sqrt{7}}}{4}$ và $\cos \alpha =\dfrac{3}{4}$ vào $P$, ta được $P=\dfrac{1}{2}.\dfrac{3}{4}+\dfrac{{\sqrt{3}}}{2}.\left({-\dfrac{{\sqrt{7}}}{4}}\right)=\dfrac{{3-\sqrt{{21}}}}{8}$.
		
	} 
\end{ex}\cham{7}

\begin{ex}%[0D6B3-4] 
	Biết $\sin a=\dfrac{5}{{13}};\cos b=\dfrac{3}{5};\,\dfrac{\pi}{2}<a<\pi;0<b<\dfrac{\pi}{2}.$ Tính $\sin \left({a+b}\right).$
	\choice
	{$\dfrac{{56}}{{65}}$}
	{$\dfrac{{63}}{{65}}$}
	{\True $-\dfrac{{33}}{{65}}$}
	{$0$}
	\loigiai{ Ta có $\cos ^2a=1-\sin ^2a=1-\left({\dfrac{5}{{13}}}\right)^2=\dfrac{{144}}{{169}}$ mà $a\in \left({\dfrac{\pi}{2};\pi}\right)\Rightarrow \cos a=-\dfrac{{12}}{{13}}.$\\
		Tương tự, ta có $\sin ^2b=1-\cos ^2b=1-\left({\dfrac{3}{5}}\right)^2=\dfrac{{16}}{{25}}$ mà $b\in \left({0;\dfrac{\pi}{2}}\right)\Rightarrow \sin b=\dfrac{4}{5}.$\\
		Khi đó $\sin \left({a+b}\right)=\sin a\cdot\cos b+\sin b.\cos a=\dfrac{5}{{13}}\cdot \dfrac{3}{5}-\dfrac{{12}}{{13}}\cdot\dfrac{4}{5}=-\dfrac{{33}}{{65}}.$
	} 
\end{ex}\cham{8}

\begin{ex}%[0D6B3-4] 
	Cho hai góc nhọn $a;b$ thoả $\cos a=\dfrac{1}{3};\cos b=\dfrac{1}{4}.$ Tính $\cos \left({a+b}\right)\cdot\cos \left({a-b}\right).$
	\choice
	{$-\dfrac{{113}}{{144}}$}
	{$-\dfrac{{115}}{{144}}$}
	{$-\dfrac{{117}}{{144}}$}
	{\True $-\dfrac{{119}}{{144}}$}
	\loigiai{  Ta có $P=\cos \left({a+b}\right)\cdot\cos \left({a-b}\right)=\left({\cos a\cdot\cos b+\sin a\cdot\sin b}\right)\left({\cos a\cdot\cos b-\sin a\cdot\sin b}\right)$\\
		$=\left({\cos a\cdot\cos b}\right)^2-\left({\sin a\cdot\sin b}\right)^2=\cos ^2a\cdot\cos ^2b-\left({1-{\cos}^2a}\right)\cdot\left({1-{\cos}^2b}\right).$
		$=\dfrac{1}{9}\cdot\dfrac{1}{{16}}-\left({1-\dfrac{1}{9}}\right)\cdot\left({1-\dfrac{1}{{16}}}\right)=-\dfrac{{119}}{{144}}.$
	}
\end{ex}\cham{9}

\begin{ex}%[0D6B3-1]%
	Cho $\tan\alpha =3$. Khi đó giá trị của $\tan \left(\alpha+\dfrac{\pi}{4}\right)$ là
	\choice
	{$\dfrac{7}{17}$}
	{$\dfrac{17}{7}$}
	{$-4$}
	{\True $-2$}
	\loigiai{
		Ta có $\tan \left(\alpha+\dfrac{\pi}{4}\right)=\dfrac{\tan\alpha+\tan \dfrac{\pi}{4}}{1-\tan\alpha\tan \dfrac{\pi}{4}}=\dfrac{3+1}{1-3\cdot 1}=-2$.
	}
\end{ex}\cham{5}

\begin{ex}%[0D6B3-1]%
	Cho $\tan\alpha=2$. Tính $\tan\left(\alpha-\dfrac{\pi}{4}\right)$.
	\choice
	{$\dfrac{2}{3}$}
	{\True $\dfrac{1}{3}$}
	{$1$}
	{$-\dfrac{1}{3}$}
	\loigiai{
		Ta có $\tan\left(\alpha-\dfrac{\pi}{4}\right)=\dfrac{\tan\alpha-\tan\dfrac{\pi}{4}}{1+\tan\alpha\tan\dfrac{\pi}{4}}=\dfrac{1}{3}$.}
\end{ex}\cham{5}

\begin{ex}%[0D6B3-2]%
	Cho $\sin \alpha=\dfrac{3}{4}$. Khi đó $ \cos 2 \alpha$ bằng
	\choice
	{ $ \dfrac{3}{8} $ }
	{\True $ \dfrac{-1}{8} $}
	{ $ \dfrac{1}{8} $ }
	{ $ \dfrac{-31}{8} $ }
	\loigiai{
		Ta có: $\cos 2 \alpha=1-2 \sin ^{2} \alpha=1-2 \cdot\left(\dfrac{3}{4}\right)^{2}=-\dfrac{1}{8}$.
	}
\end{ex}\cham{4}

\begin{ex}%[0D6B3-4] 
	Cho góc $\alpha $ thỏa mãn $\cos \alpha =\dfrac{5}{{13}}$ và $\dfrac{{3\pi}}{2}<\alpha <2\pi $. Khi đó $\tan 2\alpha $ bằng
	\choice
	{$P=-\dfrac{{120}}{{119}}$}
	{$P=-\dfrac{{119}}{{120}}$}
	{\True $P=\dfrac{{120}}{{119}}$}
	{$P=\dfrac{{119}}{{120}}$}
	\loigiai{ Ta có $P=\tan 2\alpha =\dfrac{{\sin 2\alpha}}{{\cos 2\alpha}}=\dfrac{{2\sin \alpha\cdot \cos \alpha}}{{2{\cos}^2\alpha-1}}$.\\
		Từ hệ thức $\sin ^2\alpha+\cos ^2\alpha =1$, suy ra $\sin \alpha =\pm \sqrt{{1-{\cos}^2\alpha}}=\pm \dfrac{{12}}{{13}}$.\\
		Do $\dfrac{{3\pi}}{2}<\alpha <2\pi $ nên ta chọn $\sin \alpha =-\dfrac{{12}}{{13}}$.\\
		Thay $\sin \alpha =-\dfrac{{12}}{{13}}$ và $\cos \alpha =\dfrac{5}{{13}}$ vào $P$, ta được $P=\dfrac{{120}}{{119}}$.
	} 
\end{ex}\cham{5}


\begin{ex}%[Đỗ Văn Bắc - THPT Hòa Thuận]%[0D6B3-2]%
	Nếu $\sin x+\cos x=\dfrac{1}{2}$ thì $\sin 2 x$ bằng
	\choice
	{ $ \dfrac{1}{4} $ }
	{\True $ \dfrac{-3}{4} $}
	{ $ \dfrac{1}{2} $ }
	{ $ \dfrac{-1}{2} $ }
	\loigiai{
		Ta có: $\sin x+\cos x=\dfrac{1}{2} \Rightarrow(\sin x+\cos x)^{2}=\dfrac{1}{4} \Leftrightarrow 1+\sin 2 x=\dfrac{1}{4} \Leftrightarrow \sin 2 x=-\dfrac{3}{4}$.
	}
\end{ex}\cham{5}

\begin{ex}%[0D6B3-4] 
	Cho góc $\alpha $ thỏa mãn $\sin 2\alpha =-\dfrac{4}{5}$ và $\dfrac{{3\pi}}{4}<\alpha <\pi $. Khi đó $\sin \alpha-\cos \alpha $ bằng
	\choice
	{\True $\dfrac{3}{{\sqrt{5}}}$}
	{$-\dfrac{3}{{\sqrt{5}}}$}
	{$\dfrac{{\sqrt{5}}}{3}$}
	{$-\dfrac{{\sqrt{5}}}{3}$}
	\loigiai{  Vì $\dfrac{{3\pi}}{4}<\alpha <\pi $ suy ra $\heva{& \sin \alpha >0 \\
			& \cos \alpha <0}$ nên $\sin \alpha-\cos \alpha >0$.\\
		Ta có $\left({\sin \alpha-\cos \alpha}\right)^2=1-\sin 2\alpha =1+\dfrac{4}{5}=\dfrac{9}{5}$.\\
		Suy ra $\sin \alpha-\cos \alpha =\pm \dfrac{3}{{\sqrt{5}}}$.
		Do $\sin \alpha-\cos \alpha >0$ nên $\sin \alpha-\cos \alpha =\dfrac{3}{{\sqrt{5}}}$.\\
		Vậy $P=\dfrac{3}{{\sqrt{5}}}.$
	} 
\end{ex}\cham{5}

\begin{ex}%[0D6B3-4] 
	Cho góc $\alpha $ thỏa mãn $\cos 2\alpha =-\dfrac{2}{3}$. Khi đó $\left({1+3{\sin}^2\alpha}\right)\left({1-4{\cos}^2\alpha}\right)$ bằng
	\choice
	{$12$}
	{$\dfrac{{21}}{2}$}
	{$6$}
	{\True $21$}
	\loigiai{  Ta có $P=\left({1+3\cdot\dfrac{{1-\cos 2\alpha}}{2}}\right)\left({1-4\cdot\dfrac{{1+\cos 2\alpha}}{2}}\right)=\left({\dfrac{5}{2}-\dfrac{3}{2}\cos 2\alpha}\right)\left({-1-2\cos 2\alpha}\right)$.\\
		Thay $\cos 2\alpha =-\dfrac{2}{3}$ vào $P$, ta được $P=\left({\dfrac{5}{2}+1}\right)\left({-1+\dfrac{4}{3}}\right)=\dfrac{7}{6}$.
	} 
\end{ex}\cham{5}


\begin{ex}%[0D6B3-4] 
	Cho $0<\alpha, \beta <\dfrac{\pi}{2}$ và thỏa mãn $\tan \alpha =\dfrac{1}{7}$, $\tan \beta =\dfrac{3}{4}$. Góc $\alpha+\beta $ có giá trị bằng
	\choice
	{$\dfrac{\pi}{3}$}
	{\True $\dfrac{\pi}{4}$}
	{$\dfrac{\pi}{6}$}
	{$\dfrac{\pi}{2}$}
	\loigiai{  Ta có $\tan \left({\alpha+\beta}\right)=\dfrac{{\tan \alpha+\tan \beta}}{{1-\tan \alpha\cdot\tan \beta}}=\dfrac{{\dfrac{1}{7}+\dfrac{3}{4}}}{{1-\dfrac{1}{7}\cdot\dfrac{3}{4}}}=1$ suy ra $a+b=\dfrac{\pi}{4}.$ } 
\end{ex}\cham{6}

\begin{ex}%[0D6B3-4] 
	Cho $0<x, y<\dfrac{\pi}{2}$ thỏa mãn $\cot x=\dfrac{3}{4}$, $\cot y=\dfrac{1}{7}.$ Tổng $x+y$ bằng
	\choice
	{$\dfrac{\pi}{4}$}
	{\True $\dfrac{{3\pi}}{4}$}
	{$\dfrac{\pi}{3}$}
	{$\pi$}
	\loigiai{  Ta có $\cot \left({x+y}\right)=\dfrac{{\cot x\cdot\cot y-1}}{{\cot x+\cot y}}=\dfrac{{\dfrac{3}{4}\cdot\dfrac{1}{7}-1}}{{\dfrac{3}{4}+\dfrac{1}{7}}}=-1.$\\
		Mặt khác $0<x,y<\dfrac{\pi}{2}$ suy ra $0<x+y<\pi.$ Do đó $x+y=\dfrac{{3\pi}}{4}.$
	} 
\end{ex}\cham{6}
\begin{ex}%[DCHT Toán 11 - KNTT -Phong Lê] %[1K1K2-1]
	Cho tam giác $ABC$ cân tại $A$ có $\cos B = \dfrac{5}{13}$. Tính $\sin A$.
	\choice{$\dfrac{119}{169}$}{$1$}{\True $\dfrac{120}{169}$}{$\dfrac{5}{13}$}
	\loigiai{
		Vì $0<\widehat{B}<\pi$ nên $\sin B >0$. Ta có $\sin^2B+\cos^2B = 1$, suy ra $\sin B = \sqrt{1-\left(\dfrac{5}{13}\right)^2} = \dfrac{12}{13}$.\\
		Vì tam giác $ABC$ cân tại $A$ nên $\widehat{B} = \widehat{C}$, suy ra $\sin C = \dfrac{12}{13}$ và $\cos C=\dfrac{5}{13}$.\\
		Theo tính chất tổng ba góc trong một tam giác, ta có $\widehat{A} + \widehat{B} + \widehat{C} = \pi$.\\
		Do đó
		$\sin A = \sin (\pi - A) =\sin (B+C) = \sin B\cos C + \cos B \sin C = \dfrac{12}{13}.\dfrac{5}{13} + \dfrac{5}{13}.\dfrac{12}{13} = \dfrac{120}{169}$.
	}
\end{ex}
\begin{ex}%[0D6K3-4] 
	Nếu $\tan \alpha $ và $\tan \beta $ là hai nghiệm của phương trình $x^2+px+q=0 \,\left({q\ne 1}\right)$ thì $\tan \left({\alpha+\beta}\right)$ bằng
	\choice
	{\True $\dfrac{p}{{q-1}}$}
	{$-\dfrac{p}{{q-1}}$}
	{$\dfrac{{2p}}{{1-q}}$}
	{$-\dfrac{{2p}}{{1-q}}$}
	\loigiai{  Vì $\tan \alpha,\tan \beta $ là hai nghiệm của phương trình $x^2+px+q=0$ nên theo định lí Viet, ta có$\heva{& \tan \alpha+\tan \beta =-p \\
			& \tan \alpha.\tan \beta =q}.$ Khi đó $\tan \left({\alpha+\beta}\right)=\dfrac{{\tan \alpha+\tan \beta}}{{1-\tan \alpha \tan \beta}}=\dfrac{p}{{q-1}}.$
	} 
\end{ex}\cham{6}

\begin{ex}
	Cho hai dao động điều hòa cùng phương có phương trình lần lượt là $\mathrm{x}_1=5 \cos (100 \pi \mathrm{t}+\pi)(\mathrm{cm})$ và $\mathrm{x}_2=5 \cos (100 \pi \mathrm{t}-\pi / 2)(\mathrm{cm})$. Phương trình dao động tổng hợp của hai dao động trên là
	\choice
	{$x=5 \sqrt{2} \cos \left( 100 \pi t+\dfrac{3\pi}{4}\right)\, (\mathrm{cm})$}
	{\True $x=5 \sqrt{2} \cos \left(100 \pi t-\dfrac{3\pi}{4}\right)\,(\mathrm{cm})$}
	{$x=10 \cos \left(100 \pi t-\dfrac{3\pi}{4}\right)\,(\mathrm{cm})$}
	{$x=10 \cos \left(100 \pi t+\dfrac{3\pi}{4}\right)\,(\mathrm{cm})$}
	\loigiai{
		Ta có 
		\begin{eqnarray*}
			x=x_1+x_2
			&=&5 \cos (100 \pi \mathrm{t}+\pi)+5 \cos (100 \pi \mathrm{t}-\pi / 2)\\
			&=& 5\left[\cos (100 \pi \mathrm{t}+\pi)+\cos (100 \pi \mathrm{t}-\pi / 2)\right]\\
			&=& 5 \cdot 2 cos\left(100 \pi t + \dfrac{\pi}{4} \right) \cos \dfrac{3\pi}{4}   \\
			&=& -5 \sqrt{2} \cos\left(100 \pi t + \dfrac{\pi}{4} \right)\\
			&=& 5 \sqrt{2} \cos \left(100 \pi t-\dfrac{3\pi}{4}\right).
		\end{eqnarray*}
		
	}
\end{ex}\cham{6}

\begin{ex}
	Một vật thực hiện đồng thời hai dao động điều hòa cùng phương, cùng tần số theo các phương trình: $x_1=2 \cos \left(5 \pi t+\dfrac{\pi}{2}\right)(\mathrm{cm})$ ; $x_2=2 \cos (5 \pi t)(\mathrm{cm})$. Biên độ của dao động tổng hợp của hai dao động trên là
	\choice
	{ $2$}
	{ $4$}
	{\True $2\sqrt{2}$}
	{ $\sqrt{2}$}
	\loigiai{
		Xét $x(t)=x_1+x_2=2 \cos \left(5 \pi t+\dfrac{\pi}{2}\right)+2 \cos (5 \pi t)=2\sqrt{2}\cos\left(5\pi t+\dfrac{\pi}{4} \right) $}
\end{ex}\cham{6}
\Closesolutionfile{ans}

\ind{PHẦN II.} \inden{Câu trắc nghiệm đúng sai. Học sinh trả lời từ câu 21 đến câu 25. Trong mỗi ý a), b), c), d) ở mỗi câu, học sinh chọn đúng hoặc sai.}\\
\Opensolutionfile{ans}[ans/1D1-B2-2]

\begin{ex}%[1D1B3-1]
	Cho $\sin\alpha=\dfrac{12}{13}$ và $\dfrac{\pi}{2}<\alpha<\pi$. 
	Xét tính đúng sai của các khẳng định sau:
	\choiceTF
	{\True $\cos \alpha =- \dfrac{5}{13}$}
	{\True $\cos 2\alpha=- \dfrac{119}{169}$}
	{\True $\sin\left(\alpha + \dfrac{\pi}{3}\right)=\dfrac{12 - 5\sqrt{3}}{26}$}
	{\True $\tan\left(\alpha + \dfrac{\pi}{4}\right)= \dfrac{7}{17}$}
	\loigiai{
		\begin{itemchoice}
			\itemch Vì $\dfrac{\pi}{2}<\alpha<\pi$ nên $\cos\alpha= - \sqrt{1 - \sin^2\alpha}= - \sqrt{1 - \left(\dfrac{12}{13}\right)^2}= - \dfrac{5}{13}$.
			\itemch $\cos 2\alpha=2\cos^2\alpha - 1=2\cdot\left( - \dfrac{5}{13}\right)^2 - 1= - \dfrac{119}{169}$.
			\itemch $\sin\left(\alpha + \dfrac{\pi}{3}\right)=\sin\alpha\cos\dfrac{\pi}{3} + \cos\alpha\sin\dfrac{\pi}{3}=\dfrac{12}{13}\cdot\dfrac{1}{2} + \left( - \dfrac{5}{13}\right)\cdot\dfrac{\sqrt{3}}{2}=\dfrac{12 - 5\sqrt{3}}{26}$.
			\itemch Ta có $\tan\alpha=\dfrac{\sin\alpha}{\cos\alpha}=\dfrac{\dfrac{12}{13}}{ - \dfrac{5}{13}}= - \dfrac{12}{5}$.\\
			Suy ra $\tan\left(\alpha + \dfrac{\pi}{4}\right)=\dfrac{\tan\alpha + \tan\dfrac{\pi}{4}}{1 - \tan\alpha\tan\dfrac{\pi}{4}}=\dfrac{ - \dfrac{12}{5} + 1}{1 - \left( - \dfrac{12}{5}\right)\cdot 1}=  \dfrac{7}{17}$.
		\end{itemchoice}
		}
\end{ex}\cham{21}

\begin{ex}%[1T1K3-4]
	Cho $\sin\alpha = \dfrac{3}{5}$, $\cos\beta = \dfrac{12}{13}$ và $0^\circ < \alpha, \beta < 90^\circ$. Xét tính đúng sai của các khẳng định sau:
	\choiceTF
	{$\cos \alpha =-\dfrac{4}{5}$}
	{$\sin \beta =\dfrac{25}{169}$}
	{\True $\sin\left( \alpha + \beta \right) = \dfrac{56}{65}$}
	{\True $\cos\left( \alpha - \beta \right) = \dfrac{63}{65}$}
	\loigiai{
		\begin{itemchoice}
			\itemch Vì $0^\circ < \alpha < 90^\circ$ nên $\cos\alpha > 0$. Do đó
			$$ \cos\alpha = \sqrt{1-\sin^2 \alpha} = \sqrt{1-\left( \dfrac{3}{5} \right)^2} = \dfrac{4}{5}.$$
			\itemch Vì $0^\circ < \beta < 90^\circ$ nên $\sin\beta > 0$. Do đó
			$$ \sin\beta = \sqrt{1-\cos^2 \beta} = \sqrt{1-\left( \dfrac{12}{13} \right)^2} = \dfrac{5}{13}.$$
			\itemch $ \sin\left( \alpha + \beta \right) = \sin\alpha \cos\beta + \cos\alpha \sin\beta = \dfrac{3}{5}\cdot \dfrac{12}{13} + \dfrac{4}{5}\cdot \dfrac{5}{13} = \dfrac{56}{65};$
			\itemch $ \cos\left( \alpha - \beta \right) = \cos\alpha \cos\beta + \sin\alpha \sin\beta = \dfrac{4}{5}\cdot \dfrac{12}{13} + \dfrac{3}{5}\cdot \dfrac{5}{13} = \dfrac{63}{65}.$
		\end{itemchoice}
	}
\end{ex}\cham{16}


\begin{ex}%[1D1B3-1]
	Biết $\cos 2\alpha= - \dfrac{7}{25}$ và $0<\alpha<\dfrac{\pi}{2}$.
	Xét tính đúng sai của các khẳng định sau:
	\choiceTF
	{$\cos \alpha =\dfrac{9}{25}$}
	{$\sin \alpha =\dfrac{16}{25}$}
	{$\tan \alpha =\dfrac{3}{4}$}
	{$\cot \alpha =\dfrac{4}{3}$}
	\loigiai{
		Ta có $\cos 2\alpha= - \dfrac{7}{25}$ nên $2\cos^2\alpha - 1=1 - 2\sin^2\alpha= - \dfrac{7}{25}$
		\begin{itemchoice}
			\itemch Suy ra $\cos^2\alpha=\dfrac{9}{25}$. Với $0<\alpha<\dfrac{\pi}{2}$ thì $\cos \alpha >0$ nên $\cos\alpha=\dfrac{3}{5}$.
			\itemch Suy ra $\sin^2\alpha=\dfrac{16}{25}$. Với $0<\alpha<\dfrac{\pi}{2}$ thì $\sin \alpha >0$ nên $\sin\alpha=\dfrac{4}{5}$.
			\itemch $\tan\alpha=\dfrac{\sin\alpha}{\cos\alpha}=\dfrac{\dfrac{4}{5}}{\dfrac{3}{5}}=\dfrac{4}{3}$.
			\itemch $\cot\alpha=\dfrac{1}{\tan\alpha}=\dfrac{3}{4}$.
		\end{itemchoice}
		}
\end{ex}\cham{16}

\begin{ex}%Câu 30
	\immini[thm]{Trên một mảnh giấy hình vuông $ABCD$, bác An đặt một chiếc đèn pin tại vị trí $A$ chiếu chùm sáng phân kì sang góc $C$. Bác An nhận thấy góc chiếu sáng của đèn pin giới hạn bởi hai tia $AM$ và $AN$, ở đó các điểm $M,N$ lần lượt thuộc $BC,CD$ sao cho $BM=\dfrac{1}{2}BC,DN=\dfrac{1}{3}DC$ (Hình bên).
	Xét tính đúng sai của các khẳng định sau:
		\choiceTF
		{\True $\tan \widehat{BAM}=\dfrac{1}{2}$}
		{\True $\tan \widehat{DAN}=\dfrac{1}{3}$}
		{$\widehat{BAM}+\widehat{DAN}=60^\circ $}
		{Góc chiếu sáng của đèn pin bằng $30^\circ$}
	}
	{
		\begin{tikzpicture}
			\def\a{3.5}
			\path 	(0:0) coordinate (A)
			++(0:\a) coordinate (B)
			++(90:\a) coordinate (C)
			($(C)-(B)+(A)$) coordinate (D)
			($(C)!.5!(B)$) coordinate (M)
			($(D)!1/3!(C)$) coordinate (N);
			\draw[thick] 	(A)--(B)--(C)--(D)--cycle
			(M)--(A)--(N);
			\foreach \x/\g in {A/180,B/0,C/45,D/135,M/0,N/90}
			\fill[black] 	(\x) circle (1pt)
			($(\g:3mm)+(\x)$) node {$\x$};	
	\end{tikzpicture}}
	\loigiai{
		 Giả sử hình vuông có độ dài cạnh là $1$.
		 \begin{itemchoice}
		 	\itemch Tam giác $ABM$ vuông tại $B$ có $\tan \widehat{BAM}=\dfrac{BM}{AB}=\dfrac{1}{2}$.
		 	\itemch Tam giác $ADN$ vuông tại $D$ có $\tan \widehat{DAN}=\dfrac{DN}{AD}=\dfrac{1}{3}$.
		 	\itemch $\tan \left( \widehat{BAM}+\widehat{DAN} \right)=\dfrac{\tan \widehat{BAM}+\tan \left( DAN \right)}{1-\tan \widehat{BAM}\tan \left( DAN \right)}=1\Rightarrow \widehat{BAM}+\widehat{DAN}=45^\circ$
		 	\itemch Góc chiếu sáng là $\widehat{MAN}$. Ta có
		 	$$\widehat{MAN}=90^\circ -\left( \widehat{BAM}+\widehat{DAN} \right)=45^\circ $$
		 	Vậy góc chiếu sáng của đèn pin bằng $45^\circ $.
		 \end{itemchoice}
	}
\end{ex}\cham{7}

\begin{ex}%[1D1B3-6]
	Phương trình dao động điều hòa của một vật tại thời điểm $t$ giây được cho bởi công thức $x(t)=A\cos \left(\omega t+\varphi \right)$, trong đó $x(t) $ (cm)  là li độ của vật tại thời điểm $t$ giây, $A$ là biên độ dao động $(A>0)$, $\varphi \in [-\pi; \pi]$ là pha ban đầu của dao động và $\omega$ là tần số góc. Xét hai dao động điều hòa có phương trình lần lượt là
	$$x_1(t)=3\cos\left(\dfrac{\pi}{4}t+\dfrac{\pi}{3}\right)  (\text{cm})\ \text{và} \ x_2(t)=3\cos\left(\dfrac{\pi}{4}t-\dfrac{\pi}{6}\right)  (\text{cm}). $$
	Đặt $x(t)=x_1(t)+x_2(t)$ là phương trình tổng hợp của hai dao động này.
	Xét tính đúng sai của các khẳng định sau:
	\choiceTF
	{\True Tần số góc của dao động tổng hợp là $\dfrac{\pi}{4}$}
	{Phương trình của dao động tổng hợp là $x(t)=3\cos\left(\dfrac{\pi}{4}t+\dfrac{\pi}{12}\right)$}
	{Biên độ của dao động tổng hợp trên bằng $3$}
	{\True Pha ban đầu của dao động tổng hợp trên là $\dfrac{\pi}{12}$ }
	\loigiai{
		\begin{itemchoice}
			\itemch Tần số góc $\omega=\dfrac{\pi}{4}$.
			\itemch Ta có 
			\begin{eqnarray*}
				x(t)&=&x_1(t)+x_2(t)\\
				&=&3\cos\left(\dfrac{\pi}{4}t+\dfrac{\pi}{3}\right) +3\cos\left(\dfrac{\pi}{4}t-\dfrac{\pi}{6}\right)\\
				&=&3.2\cos \dfrac{\left(\dfrac{\pi}{4}t+\dfrac{\pi}{3}\right)+\left(\dfrac{\pi}{4}t-\dfrac{\pi}{6}\right)}{2}\cos \dfrac{\left(\dfrac{\pi}{4}t+\dfrac{\pi}{3}\right)-\left(\dfrac{\pi}{4}t-\dfrac{\pi}{6}\right)}{2}\\
				&=& 6 \cos \dfrac{\dfrac{\pi}{2}t+\dfrac{\pi}{6}}{2} \cos \dfrac{\tfrac{\pi}{2}}{2}\\
				&=&3\sqrt{2}\cos\left(\dfrac{\pi}{4}t+\dfrac{\pi}{12}\right).
			\end{eqnarray*}
			Vậy phương trình của dao động tổng hợp là $x_(t)=3\sqrt{2}\cos\left(\dfrac{\pi}{4}t+\dfrac{\pi}{12}\right).$
			\itemch Căn cứ vào phương trình $x(t)$, ta có biên độ dao động là $3\sqrt{2}$.
			\itemch Căn cứ vào phương trình $x(t)$, ta có pha ban đầu là $\dfrac{\pi}{12}$.
		\end{itemchoice}
	}
\end{ex}\cham{10}

\Closesolutionfile{ans}

\ind{PHẦN III.} \inden{Câu trắc nghiệm trả lời ngắn. Học sinh trả lời từ câu 26 đến câu 30 vào ô kết quả.}\\
\Opensolutionfile{ans}[ans/1D1-B2-3]
\begin{ex}%[1D1V3-4]
Cho $\cos x=\dfrac{2}{3}$, $\sin y=\dfrac{1}{5}$. Tính giá trị của biểu thức $\cos (x - y)\cos (x + y)$  (kết quả làm tròn đến chữ số hàng phần chục).
\shortans[0]{$0,4$}
\loigiai{
Ta có
\begin{align*}
\cos (x - y)\cos (x + y) &= \dfrac{1}{2}\left[\cos (2x) +\cos (-2y)\right] \\
&=\dfrac{1}{2}\left[2\cos^2 x-1 + 1-2\sin^2 y\right] \\
&=\cos^2 x-\sin^2 y \\
&=\left(\dfrac{2}{3}\right)^2-\left(\dfrac{1}{5}\right)^2=\dfrac{91}{225}\approx 0{,}4.
\end{align*}
}

\end{ex}
\begin{ex}%Câu 15
	Cho hai góc $a$ và $b$ với $\tan a=\dfrac{1}{7}$ và $\tan b=\dfrac{3}{4}$. Khi đó, $\tan \left( a+b \right)$ bằng bao nhiêu?\\
	\shortans[oly]{1}
	\loigiai{
		Ta có $\tan \left( a+b \right)=\dfrac{\tan a+\tan b}{1-\tan a\tan b}$.\\
		Mà $\tan a=\dfrac{1}{7}$ và $\tan b=\dfrac{3}{4}$ nên $\tan \left( a+b \right)=1$.}
\end{ex}\cham{4}

\begin{ex}
	\immini[thm]{Cho hình vuông $ABCD$ như hình vẽ bên. Tính $\tan x$ (kết quả làm tròn đến hàng phần trăm).\\
	\shortans[oly]{$0{,}41$}}{
\begin{tikzpicture}[scale=0.75, font=\footnotesize,>=stealth]
	\path
	%	Vẽ mp
	(0,0) coordinate (A)
	(3,0) coordinate (B)
	(3,3) coordinate (C)
	(0,3) coordinate (D)
	(3,1.2) coordinate (E)
	;
	\draw (A)--(B)--(E)node[pos=.5,right]{$5$}--(C)node[pos=.5,right]{$7$}--(D)--(A) (C)--(A)--(E);
	\draw pic["\scriptsize$x$",draw,angle eccentricity=1.6,angle radius=0.5cm]{angle=E--A--C};
	\foreach \x/\g in {A/180,B/0,C/90,D/90,E/0}\draw[fill=black] (\x) circle (.05) +(\g:.5)node{\footnotesize$\x$};
\end{tikzpicture}}
\loigiai{
	Ta có
		\begin{eqnarray*}
		&&\tan\left(45^\circ-x \right)=\tan \widehat{EAB}=\dfrac{5}{12}\\
		&\Leftrightarrow& \dfrac{1-\tan x}{1+\tan x}=\dfrac{5}{12}\\
		&\Leftrightarrow& \tan x = \dfrac{7}{17}\approx 0{,}41.
	\end{eqnarray*}
}
\end{ex}\cham{5}

% \begin{ex}
% 	Biết $\dfrac{\sin \alpha+\sin 2 \alpha}{1+\cos \alpha+\cos 2 \alpha}=n \cdot \tan \alpha$, $n \in \mathbb{N}$. Giá trị $n$ bằng bao nhiêu?\\
% 	\shortans[oly]{1}
% 	\loigiai{
% 		Xét $$\dfrac{\sin \alpha+\sin 2 \alpha}{1+\cos \alpha+\cos 2 \alpha}=\dfrac{\sin \alpha+2\sin  \alpha \cos\alpha}{1+\cos \alpha+2\cos^2 \alpha-1}=\dfrac{\sin \alpha\left(1+2\cos \alpha \right)}{\cos \alpha\left(1+2\cos \alpha \right)}=\dfrac{\sin\alpha}{\sin\alpha}=\tan \alpha.$$
% 		Suy ra $n=1$.
% 	}
% \end{ex}\cham{5}
\begin{ex}%[Hoàng Thanh Phương, 24-25 BG K11]%[1D1V3-3]
Biết $\sin^4x=a+b\cos 2x+c \cdot \cos 4x$, với $a,b,c \in \mathbb{Q}$. Tính tổng $S=a+b+c$.
\shortans{$0$}
\loigiai{
Ta có \allowdisplaybreaks
\begin{eqnarray*}
& \sin^4x & =\left(\dfrac{1-\cos 2x}{2}\right)^2=\dfrac{1}{4}\left(1-2\cos 2x+\cos^22x\right)\\
& & =\dfrac{1}{4}\left(1-2\cos 2x+\dfrac{1+\cos 4x}{2}\right) =\dfrac{3}{8}-\dfrac{1}{2}\cos 2x+\dfrac{1}{8}\cos 4x.
\end{eqnarray*}
Vậy $a=\dfrac{3}{8}$, $b=-\dfrac{1}{2}$, $c=\dfrac{1}{8}$. Nên $S=a+b+c=0$.}
\end{ex}

\begin{ex}
	\immini[thm]{Cho tam giác $ABC$ như hình vẽ. Biết $\widehat{B}=2\widehat{C}=2\alpha$, $AB=2$, $AC=3$. Tính độ dài cạnh $BC$.\\
		\shortans[oly]{$2,5$}}{
		\begin{tikzpicture}[scale=1, font=\footnotesize,>=stealth]
			\path
			%	Vẽ mp
			(1,2) coordinate (A)
			(0,0) coordinate (B)
			(4.4,0) coordinate (C)
			;
			\draw (A)--(B)node[pos=.5,left]{$2$}--(C)node[pos=.5,below]{$x$}--(A)node[pos=.5,right]{$3$};
			\draw pic["\scriptsize$2\alpha$",draw,angle eccentricity=1.6,angle radius=0.5cm]{angle=C--B--A};
			\draw pic["\scriptsize$\alpha$",draw,angle eccentricity=1.6,angle radius=0.5cm]{angle=A--C--B};
			\foreach \x/\g in {A/90,B/180,C/0}\draw[fill=black] (\x) circle (.05) +(\g:.5)node{\footnotesize$\x$};
	\end{tikzpicture}}
\loigiai{

Đặt $BC = x$ ($x > 0$).
Áp dụng định lí sin cho tam giác $ABC$, ta có:
$$ \dfrac{AC}{\sin B} = \dfrac{AB}{\sin C} \Leftrightarrow \dfrac{3}{\sin(2\alpha)} = \dfrac{2}{\sin \alpha} $$
$$ \Leftrightarrow \dfrac{3}{2\sin\alpha\cos\alpha} = \dfrac{2}{\sin\alpha} \Rightarrow \cos\alpha = \dfrac{3}{4}. $$

Áp dụng định lí côsin cho tam giác $ABC$, ta có:
$$ AB^2 = BC^2 + AC^2 - 2 \cdot BC \cdot AC \cdot \cos C $$
$$ \Leftrightarrow 2^2 = x^2 + 3^2 - 2 \cdot x \cdot 3 \cdot \dfrac{3}{4} $$
$$ \Leftrightarrow x^2 - \dfrac{9}{2}x + 5 = 0 \Leftrightarrow \hoac{x &= \dfrac{5}{2} \\ x &= 2.} $$

Kiểm tra nghiệm:
\begin{itemize}
	\item Nếu $x = 2$, tam giác $ABC$ có $AB = BC = 2$ nên cân tại $B$. Khi đó $\widehat{A} = \widehat{C} = \alpha \Rightarrow \widehat{B} = 180^\circ - 2\alpha$.
	Mà giả thiết $\widehat{B} = 2\alpha$ nên $180^\circ - 2\alpha = 2\alpha \Leftrightarrow \alpha = 45^\circ$. Suy ra $\cos \alpha = \dfrac{\sqrt{2}}{2}$ (mâu thuẫn với $\cos\alpha = \dfrac{3}{4}$). Vậy loại $x = 2$.
	\item Nếu $x = \dfrac{5}{2} = 2{,}5$, áp dụng lại định lí côsin ta được $\cos B = \dfrac{BC^2 + AB^2 - AC^2}{2 \cdot BC \cdot AB} = \dfrac{2{,}5^2 + 2^2 - 3^2}{2 \cdot 2{,}5 \cdot 2} = \dfrac{1}{8}$. 
	Lại có $\cos(2C) = 2\cos^2\alpha - 1 = 2\left(\dfrac{3}{4}\right)^2 - 1 = \dfrac{1}{8}$. Suy ra $\cos B = \cos(2C)$, thỏa mãn giả thiết $\widehat{B} = 2\widehat{C}$.
\end{itemize}

Vậy độ dài cạnh $BC$ là $2{,}5$.}
\end{ex}\cham{5}

\begin{ex}
	\immini[thm]{Trên nóc một tòa nhà thẳng đứng có một cột ăng-ten cao $5 \mathrm{~m}$. Từ một vị trí quan sát $A$ cao $7 \mathrm{~m}$ so với mặt đất và cách tòa nhà 14 m có thể nhìn thấy đỉnh $B$ và chân $C$ của cột ăng-ten với góc $\widehat{BAC}=10^\circ$. Tính chiều cao của tòa nhà (tính bằng mét và kết quả làm tròn đến hàng đơn vị).\\
		\shortans[oly]{19}
}{
	\begin{tikzpicture}[scale=1, font=\footnotesize,>=stealth]
		\path 
		(0,0) coordinate (A)
		(4,3) coordinate (C)
		(4,4) coordinate (B)
		(4,-1) coordinate (O) 
		(0,-1) coordinate (L) 
		(4,0) coordinate (R) 
		;
		\draw[dashed] (A)--(B) (A)--(C);
		\draw[thick](B)--(C);
		\draw[fill=cyan] (C)--(6,3)--(6,-1)--(4,-1)--(C);
		\foreach \i in {2,3,4,5}
		\draw[fill=white] ({4+1/3},{-1+(2/3)*(\i-1)})--({4+1/3},{-1+(2/3)*(\i-1)+0.5})--({4+2/3},{-1+(2/3)*(\i-1)+0.5})--({4+2/3},{-1+(2/3)*(\i-1)})--cycle 
		({5+1/3},{-1+(2/3)*(\i-1)})--({5+1/3},{-1+(2/3)*(\i-1)+0.5})--({5+2/3},{-1+(2/3)*(\i-1)+0.5})--({5+2/3},{-1+(2/3)*(\i-1)})--cycle;
		\draw[dashed]   (A)--(R)node[pos=.5,below]{$14$ m}; 
		\draw [<->](A)--(L);
		\draw ($(A)!0.5!(L)$) node[left]{$7 \mathrm{~m}$} ($(B)!0.5!(C)$) node[right]{\scriptsize$5 \mathrm{~m}$} ;
		\draw pic[draw=black,double, angle eccentricity=1.4, angle radius=0.6cm]{angle=C--A--B};
		\foreach \x/\g in {A/180,B/140,C/140} 
		\fill[black] (\x) circle (1pt)+(\g:3mm) node {$\x$};
	%	\clip (-0.5,{-1-0.25}) rectangle (6.6,-1) ;
		\draw[fill=gray!30] (-0.5,-1)--(6.6,-1)--(6.6,{-1-0.25})--(-0.5,{-1-0.25})--cycle;
		\foreach \i in {1,2,...,35}
		\draw ({-0.5+0.2*(\i)},-1)--($({-0.5+0.2*(\i)},-1)+(-120:0.3)$) ;
		\node[above] at (0.7,1) {\scriptsize$10^\circ$};
		\draw[->] (0.7,1)--(0.5,0.45);
\end{tikzpicture}}
\loigiai{
	\immini[thm]{
	Đặt $CR=h>0$ và $\widehat{CAR}=\alpha$. Ta có
	\begin{itemize}
		\item [$\bullet$] $\tan\alpha=\dfrac{h}{14}$.
		\item [$\bullet$] $\tan\left( \alpha+10^\circ\right) =\dfrac{h+5}{14}$.
	\end{itemize}
	$\tan\left( \alpha+10^\circ\right) =\dfrac{h+5}{14} \Leftrightarrow \dfrac{\tan\alpha + \tan10^\circ}{1-\tan \alpha.\tan10^\circ}=\dfrac{h+5}{14}$.\\
	$\Leftrightarrow \dfrac{\dfrac{h}{14} + \tan10^\circ}{1-\dfrac{h}{14}.\tan10^\circ}=\dfrac{h+5}{14}$. Giải phương trình này, ta tìm được $h \approx 12$. Suy ra, chiều cao ngôi nhà là $h+7=19$ (m).
	}{
	\begin{tikzpicture}[scale=1, font=\footnotesize,>=stealth]
		\path 
		(0,0) coordinate (A)
		(4,3) coordinate (C)
		(4,4) coordinate (B)
		(4,-1) coordinate (O) 
		(0,-1) coordinate (L) 
		(4,0) coordinate (R) 
		;
		\draw[dashed] (A)--(B) (A)--(C);
		\draw[thick](B)--(C);
		\draw[fill=cyan] (C)--(6,3)--(6,-1)--(4,-1)--(C);
		\foreach \i in {2,3,4,5}
		\draw[fill=white] ({4+1/3},{-1+(2/3)*(\i-1)})--({4+1/3},{-1+(2/3)*(\i-1)+0.5})--({4+2/3},{-1+(2/3)*(\i-1)+0.5})--({4+2/3},{-1+(2/3)*(\i-1)})--cycle 
		({5+1/3},{-1+(2/3)*(\i-1)})--({5+1/3},{-1+(2/3)*(\i-1)+0.5})--({5+2/3},{-1+(2/3)*(\i-1)+0.5})--({5+2/3},{-1+(2/3)*(\i-1)})--cycle;
		\draw[dashed]   (A)--(R)node[pos=.5,below]{$14$ m}; 
		\draw [<->](A)--(L);
		\draw ($(A)!0.5!(L)$) node[left]{$7 \mathrm{~m}$} ($(B)!0.5!(C)$) node[right]{\scriptsize$5 \mathrm{~m}$} ;
		\draw pic[draw=black,double, angle eccentricity=1.4, angle radius=0.6cm]{angle=C--A--B};
		\foreach \x/\g in {A/180,B/140,C/140,R/-120} 
		\fill[black] (\x) circle (1pt)+(\g:3mm) node {$\x$};
		%	\clip (-0.5,{-1-0.25}) rectangle (6.6,-1) ;
		\draw[fill=gray!30] (-0.5,-1)--(6.6,-1)--(6.6,{-1-0.25})--(-0.5,{-1-0.25})--cycle;
		\foreach \i in {1,2,...,35}
		\draw ({-0.5+0.2*(\i)},-1)--($({-0.5+0.2*(\i)},-1)+(-120:0.3)$) ;
		\node[above] at (0.7,1) {\scriptsize$10^\circ$};
		\draw[->] (0.7,1)--(0.5,0.45);
\end{tikzpicture}}
	
}
\end{ex}\cham{25}
\Closesolutionfile{ans}


=======
\subsection{BÀI TẬP TRẮC NGHIỆM}

\ind{PHẦN I.} \inden{Câu trắc nghiệm nhiều phương án lựa chọn. Học sinh trả lời từ câu 1 đến câu 20. Mỗi câu hỏi học sinh chỉ chọn một phương án.}\\
\setcounter{ex}{0}
\Opensolutionfile{ans}[ans/1D1-B2-1]
\begin{ex}%[0D6Y3-8]
	Khẳng định nào sau đây là khẳng định \textbf{sai}?
	\choice
	{$\cos \left({a-b}\right)=\sin a\sin b+\cos a\cos b$}
	{\True $\cos \left({a+b}\right)=\sin a\sin b-\cos a\cos b$}
	{$\sin \left({a-b}\right)=\sin a\cos b-\cos a\sin b$}
	{$\sin \left({a+b}\right)=\sin a\cos b+\cos a\sin b$}
	\loigiai{Ta có $\cos \left({a+b}\right)=\cos a\cos b-\sin a\sin b$. } 
\end{ex}\cham{2}

\begin{ex}%[0D6B3-8]
	Khẳng định nào sau đây là khẳng định đúng?
	\choice
	{$\sin a+\cos a=\sqrt{2}\sin \left({a-\dfrac{\pi}{4}}\right)$}
	{\True $\sin a+\cos a=\sqrt{2}\sin \left({a+\dfrac{\pi}{4}}\right)$}
	{$\sin a+\cos a=-\sqrt{2}\sin \left({a-\dfrac{\pi}{4}}\right)$}
	{$\sin a+\cos a=-\sqrt{2}\sin \left({a+\dfrac{\pi}{4}}\right)$}
	\loigiai{  } 
\end{ex}\cham{2}

\begin{ex}%[0D6B3-1]%
	Cho các góc lượng giác $a,b$ và $T=\cos (a+b)\cos (a-b)-\sin (a+b)\sin (a-b)$. Mệnh đề sau đây đúng?
	\choice
	{$T=\sin 2b$}
	{$T=\sin 2a$}
	{\True $T=\cos 2a$}
	{$T=\cos 2b$}
	\loigiai{
		Ta có $T=\cos (a+b)\cos (a-b)-\sin (a+b)\sin (a-b)=\cos [(a+b)+(a-b)]=\cos 2a$.}
\end{ex}\cham{6}

\begin{ex}%[0D6B3-1]%
	Cho $\sin \alpha=\dfrac{1}{\sqrt{3}}$ với $0<\alpha<\dfrac{\pi}{2}$. Khi đó, giá trị của $\sin \left(\alpha+\dfrac{\pi}{3}\right)$ bằng
	\choice
	{$\dfrac{\sqrt{3}}{3}-\dfrac{1}{2}$}
	{\True $\dfrac{\sqrt{3}}{6}+\dfrac{\sqrt{2}}{2}$}
	{$\sqrt{6}-\dfrac{1}{2}$}
	{$\dfrac{\sqrt{3}}{6}-\dfrac{\sqrt{2}}{2}$}
	\loigiai{
		Với $0<\alpha<\dfrac{\pi}{2}$ thì $\cos\alpha >0$ nên $\cos \alpha=\sqrt{1-\sin^2 \alpha}=\dfrac{\sqrt{6}}{3}$.\\
		Ta có $\sin \left(\alpha+\dfrac{\pi}{3}\right)=\sin \alpha\cdot \cos \dfrac{\pi}{3}+\cos \alpha\cdot \sin \dfrac{\pi}{3}=\dfrac{1}{\sqrt{3}}\cdot \dfrac{1}{2}+\dfrac{\sqrt{6}}{3}\cdot \dfrac{\sqrt{3}}{2}=\dfrac{\sqrt{3}}{6}+\dfrac{\sqrt{2}}{2}$.
	}
\end{ex}\cham{6}


\begin{ex}%[0D6B3-4] 
	Biết $\sin \alpha=-\dfrac{3}{5}$ và $\pi <\alpha <\dfrac{{3\pi}}{2}$. Tính $\sin \left({\alpha+\dfrac{\pi}{6}}\right).$
	\choice
	{$-\dfrac{3}{5}$}
	{$\dfrac{3}{5}$}
	{\True $\dfrac{{-4-3\sqrt{3}}}{{10}}$}
	{$\dfrac{{4-3\sqrt{3}}}{{10}}$}
	\loigiai{
		Từ hệ thức $\sin ^2\alpha+\cos ^2\alpha =1$, suy ra $\cos \alpha =\pm \sqrt{{1-{\sin}^2\alpha}}=\pm \dfrac{4}{5}$.
		Do $\pi <\alpha <\dfrac{{3\pi}}{2}$ nên ta chọn $\cos \alpha =-\dfrac{4}{5}$.\\
		Suy ra $P=\sin \left({\alpha+\dfrac{\pi}{6}}\right)=\dfrac{{\sqrt{3}}}{2}\sin \alpha+\dfrac{1}{2}\cos \alpha =\dfrac{{\sqrt{3}}}{2}\left({-\dfrac{3}{5}}\right)+\dfrac{1}{2}\left({-\dfrac{4}{5}}\right)=\dfrac{{-4-3\sqrt{3}}}{{10}}$.
	} 
\end{ex}\cham{7}

\begin{ex}%[0D6B3-4] 
	Cho góc $\alpha $ thỏa mãn $\cos \alpha =\dfrac{3}{4}$ và $\dfrac{{3\pi}}{2}<\alpha <2\pi $. Tính $\cos \left({\dfrac{\pi}{3}-\alpha}\right).$
	\choice
	{$\dfrac{{3+\sqrt{{21}}}}{8}$}
	{\True $\dfrac{{3-\sqrt{{21}}}}{8}$}
	{$\dfrac{{3\sqrt{3}+\sqrt{7}}}{8}$}
	{$\dfrac{{3\sqrt{3}-\sqrt{7}}}{8}$}
	\loigiai{ Ta có $P=\cos \left({\dfrac{\pi}{3}-\alpha}\right)=\cos \dfrac{\pi}{3}\cos \alpha+\sin \dfrac{\pi}{3}\sin \alpha =\dfrac{1}{2}\cos \alpha+\dfrac{{\sqrt{3}}}{2}\sin \alpha $.\\
		Từ hệ thức $\sin ^2\alpha+\cos ^2\alpha =1$, suy ra $\sin \alpha =\pm \sqrt{{1-{\cos}^2\alpha}}=\pm \dfrac{{\sqrt{7}}}{4}$.\\
		Do $\dfrac{{3\pi}}{2}<\alpha <2\pi $ nên ta chọn $\sin \alpha =-\dfrac{{\sqrt{7}}}{4}$.\\
		Thay $\sin \alpha =-\dfrac{{\sqrt{7}}}{4}$ và $\cos \alpha =\dfrac{3}{4}$ vào $P$, ta được $P=\dfrac{1}{2}.\dfrac{3}{4}+\dfrac{{\sqrt{3}}}{2}.\left({-\dfrac{{\sqrt{7}}}{4}}\right)=\dfrac{{3-\sqrt{{21}}}}{8}$.
		
	} 
\end{ex}\cham{7}

\begin{ex}%[0D6B3-4] 
	Biết $\sin a=\dfrac{5}{{13}};\cos b=\dfrac{3}{5};\,\dfrac{\pi}{2}<a<\pi;0<b<\dfrac{\pi}{2}.$ Tính $\sin \left({a+b}\right).$
	\choice
	{$\dfrac{{56}}{{65}}$}
	{$\dfrac{{63}}{{65}}$}
	{\True $-\dfrac{{33}}{{65}}$}
	{$0$}
	\loigiai{ Ta có $\cos ^2a=1-\sin ^2a=1-\left({\dfrac{5}{{13}}}\right)^2=\dfrac{{144}}{{169}}$ mà $a\in \left({\dfrac{\pi}{2};\pi}\right)\Rightarrow \cos a=-\dfrac{{12}}{{13}}.$\\
		Tương tự, ta có $\sin ^2b=1-\cos ^2b=1-\left({\dfrac{3}{5}}\right)^2=\dfrac{{16}}{{25}}$ mà $b\in \left({0;\dfrac{\pi}{2}}\right)\Rightarrow \sin b=\dfrac{4}{5}.$\\
		Khi đó $\sin \left({a+b}\right)=\sin a\cdot\cos b+\sin b.\cos a=\dfrac{5}{{13}}\cdot \dfrac{3}{5}-\dfrac{{12}}{{13}}\cdot\dfrac{4}{5}=-\dfrac{{33}}{{65}}.$
	} 
\end{ex}\cham{8}

\begin{ex}%[0D6B3-4] 
	Cho hai góc nhọn $a;b$ thoả $\cos a=\dfrac{1}{3};\cos b=\dfrac{1}{4}.$ Tính $\cos \left({a+b}\right)\cdot\cos \left({a-b}\right).$
	\choice
	{$-\dfrac{{113}}{{144}}$}
	{$-\dfrac{{115}}{{144}}$}
	{$-\dfrac{{117}}{{144}}$}
	{\True $-\dfrac{{119}}{{144}}$}
	\loigiai{  Ta có $P=\cos \left({a+b}\right)\cdot\cos \left({a-b}\right)=\left({\cos a\cdot\cos b+\sin a\cdot\sin b}\right)\left({\cos a\cdot\cos b-\sin a\cdot\sin b}\right)$\\
		$=\left({\cos a\cdot\cos b}\right)^2-\left({\sin a\cdot\sin b}\right)^2=\cos ^2a\cdot\cos ^2b-\left({1-{\cos}^2a}\right)\cdot\left({1-{\cos}^2b}\right).$
		$=\dfrac{1}{9}\cdot\dfrac{1}{{16}}-\left({1-\dfrac{1}{9}}\right)\cdot\left({1-\dfrac{1}{{16}}}\right)=-\dfrac{{119}}{{144}}.$
	}
\end{ex}\cham{9}

\begin{ex}%[0D6B3-1]%
	Cho $\tan\alpha =3$. Khi đó giá trị của $\tan \left(\alpha+\dfrac{\pi}{4}\right)$ là
	\choice
	{$\dfrac{7}{17}$}
	{$\dfrac{17}{7}$}
	{$-4$}
	{\True $-2$}
	\loigiai{
		Ta có $\tan \left(\alpha+\dfrac{\pi}{4}\right)=\dfrac{\tan\alpha+\tan \dfrac{\pi}{4}}{1-\tan\alpha\tan \dfrac{\pi}{4}}=\dfrac{3+1}{1-3\cdot 1}=-2$.
	}
\end{ex}\cham{5}

\begin{ex}%[0D6B3-1]%
	Cho $\tan\alpha=2$. Tính $\tan\left(\alpha-\dfrac{\pi}{4}\right)$.
	\choice
	{$\dfrac{2}{3}$}
	{\True $\dfrac{1}{3}$}
	{$1$}
	{$-\dfrac{1}{3}$}
	\loigiai{
		Ta có $\tan\left(\alpha-\dfrac{\pi}{4}\right)=\dfrac{\tan\alpha-\tan\dfrac{\pi}{4}}{1+\tan\alpha\tan\dfrac{\pi}{4}}=\dfrac{1}{3}$.}
\end{ex}\cham{5}

\begin{ex}%[0D6B3-2]%
	Cho $\sin \alpha=\dfrac{3}{4}$. Khi đó $ \cos 2 \alpha$ bằng
	\choice
	{ $ \dfrac{3}{8} $ }
	{\True $ \dfrac{-1}{8} $}
	{ $ \dfrac{1}{8} $ }
	{ $ \dfrac{-31}{8} $ }
	\loigiai{
		Ta có: $\cos 2 \alpha=1-2 \sin ^{2} \alpha=1-2 \cdot\left(\dfrac{3}{4}\right)^{2}=-\dfrac{1}{8}$.
	}
\end{ex}\cham{4}

\begin{ex}%[0D6B3-4] 
	Cho góc $\alpha $ thỏa mãn $\cos \alpha =\dfrac{5}{{13}}$ và $\dfrac{{3\pi}}{2}<\alpha <2\pi $. Khi đó $\tan 2\alpha $ bằng
	\choice
	{$P=-\dfrac{{120}}{{119}}$}
	{$P=-\dfrac{{119}}{{120}}$}
	{\True $P=\dfrac{{120}}{{119}}$}
	{$P=\dfrac{{119}}{{120}}$}
	\loigiai{ Ta có $P=\tan 2\alpha =\dfrac{{\sin 2\alpha}}{{\cos 2\alpha}}=\dfrac{{2\sin \alpha\cdot \cos \alpha}}{{2{\cos}^2\alpha-1}}$.\\
		Từ hệ thức $\sin ^2\alpha+\cos ^2\alpha =1$, suy ra $\sin \alpha =\pm \sqrt{{1-{\cos}^2\alpha}}=\pm \dfrac{{12}}{{13}}$.\\
		Do $\dfrac{{3\pi}}{2}<\alpha <2\pi $ nên ta chọn $\sin \alpha =-\dfrac{{12}}{{13}}$.\\
		Thay $\sin \alpha =-\dfrac{{12}}{{13}}$ và $\cos \alpha =\dfrac{5}{{13}}$ vào $P$, ta được $P=\dfrac{{120}}{{119}}$.
	} 
\end{ex}\cham{5}


\begin{ex}%[Đỗ Văn Bắc - THPT Hòa Thuận]%[0D6B3-2]%
	Nếu $\sin x+\cos x=\dfrac{1}{2}$ thì $\sin 2 x$ bằng
	\choice
	{ $ \dfrac{1}{4} $ }
	{\True $ \dfrac{-3}{4} $}
	{ $ \dfrac{1}{2} $ }
	{ $ \dfrac{-1}{2} $ }
	\loigiai{
		Ta có: $\sin x+\cos x=\dfrac{1}{2} \Rightarrow(\sin x+\cos x)^{2}=\dfrac{1}{4} \Leftrightarrow 1+\sin 2 x=\dfrac{1}{4} \Leftrightarrow \sin 2 x=-\dfrac{3}{4}$.
	}
\end{ex}\cham{5}

\begin{ex}%[0D6B3-4] 
	Cho góc $\alpha $ thỏa mãn $\sin 2\alpha =-\dfrac{4}{5}$ và $\dfrac{{3\pi}}{4}<\alpha <\pi $. Khi đó $\sin \alpha-\cos \alpha $ bằng
	\choice
	{\True $\dfrac{3}{{\sqrt{5}}}$}
	{$-\dfrac{3}{{\sqrt{5}}}$}
	{$\dfrac{{\sqrt{5}}}{3}$}
	{$-\dfrac{{\sqrt{5}}}{3}$}
	\loigiai{  Vì $\dfrac{{3\pi}}{4}<\alpha <\pi $ suy ra $\heva{& \sin \alpha >0 \\
			& \cos \alpha <0}$ nên $\sin \alpha-\cos \alpha >0$.\\
		Ta có $\left({\sin \alpha-\cos \alpha}\right)^2=1-\sin 2\alpha =1+\dfrac{4}{5}=\dfrac{9}{5}$.\\
		Suy ra $\sin \alpha-\cos \alpha =\pm \dfrac{3}{{\sqrt{5}}}$.
		Do $\sin \alpha-\cos \alpha >0$ nên $\sin \alpha-\cos \alpha =\dfrac{3}{{\sqrt{5}}}$.\\
		Vậy $P=\dfrac{3}{{\sqrt{5}}}.$
	} 
\end{ex}\cham{5}

\begin{ex}%[0D6B3-4] 
	Cho góc $\alpha $ thỏa mãn $\cos 2\alpha =-\dfrac{2}{3}$. Khi đó $\left({1+3{\sin}^2\alpha}\right)\left({1-4{\cos}^2\alpha}\right)$ bằng
	\choice
	{$12$}
	{$\dfrac{{21}}{2}$}
	{$6$}
	{\True $21$}
	\loigiai{  Ta có $P=\left({1+3\cdot\dfrac{{1-\cos 2\alpha}}{2}}\right)\left({1-4\cdot\dfrac{{1+\cos 2\alpha}}{2}}\right)=\left({\dfrac{5}{2}-\dfrac{3}{2}\cos 2\alpha}\right)\left({-1-2\cos 2\alpha}\right)$.\\
		Thay $\cos 2\alpha =-\dfrac{2}{3}$ vào $P$, ta được $P=\left({\dfrac{5}{2}+1}\right)\left({-1+\dfrac{4}{3}}\right)=\dfrac{7}{6}$.
	} 
\end{ex}\cham{5}


\begin{ex}%[0D6B3-4] 
	Cho $0<\alpha, \beta <\dfrac{\pi}{2}$ và thỏa mãn $\tan \alpha =\dfrac{1}{7}$, $\tan \beta =\dfrac{3}{4}$. Góc $\alpha+\beta $ có giá trị bằng
	\choice
	{$\dfrac{\pi}{3}$}
	{\True $\dfrac{\pi}{4}$}
	{$\dfrac{\pi}{6}$}
	{$\dfrac{\pi}{2}$}
	\loigiai{  Ta có $\tan \left({\alpha+\beta}\right)=\dfrac{{\tan \alpha+\tan \beta}}{{1-\tan \alpha\cdot\tan \beta}}=\dfrac{{\dfrac{1}{7}+\dfrac{3}{4}}}{{1-\dfrac{1}{7}\cdot\dfrac{3}{4}}}=1$ suy ra $a+b=\dfrac{\pi}{4}.$ } 
\end{ex}\cham{6}

\begin{ex}%[0D6B3-4] 
	Cho $0<x, y<\dfrac{\pi}{2}$ thỏa mãn $\cot x=\dfrac{3}{4}$, $\cot y=\dfrac{1}{7}.$ Tổng $x+y$ bằng
	\choice
	{$\dfrac{\pi}{4}$}
	{\True $\dfrac{{3\pi}}{4}$}
	{$\dfrac{\pi}{3}$}
	{$\pi$}
	\loigiai{  Ta có $\cot \left({x+y}\right)=\dfrac{{\cot x\cdot\cot y-1}}{{\cot x+\cot y}}=\dfrac{{\dfrac{3}{4}\cdot\dfrac{1}{7}-1}}{{\dfrac{3}{4}+\dfrac{1}{7}}}=-1.$\\
		Mặt khác $0<x,y<\dfrac{\pi}{2}$ suy ra $0<x+y<\pi.$ Do đó $x+y=\dfrac{{3\pi}}{4}.$
	} 
\end{ex}\cham{6}
\begin{ex}%[0D6K3-4] 
	Nếu $\tan \alpha $ và $\tan \beta $ là hai nghiệm của phương trình $x^2+px+q=0 \,\left({q\ne 1}\right)$ thì $\tan \left({\alpha+\beta}\right)$ bằng
	\choice
	{\True $\dfrac{p}{{q-1}}$}
	{$-\dfrac{p}{{q-1}}$}
	{$\dfrac{{2p}}{{1-q}}$}
	{$-\dfrac{{2p}}{{1-q}}$}
	\loigiai{  Vì $\tan \alpha,\tan \beta $ là hai nghiệm của phương trình $x^2+px+q=0$ nên theo định lí Viet, ta có$\heva{& \tan \alpha+\tan \beta =-p \\
			& \tan \alpha.\tan \beta =q}.$ Khi đó $\tan \left({\alpha+\beta}\right)=\dfrac{{\tan \alpha+\tan \beta}}{{1-\tan \alpha \tan \beta}}=\dfrac{p}{{q-1}}.$
	} 
\end{ex}\cham{6}

\begin{ex}
	Cho hai dao động điều hòa cùng phương có phương trình lần lượt là $\mathrm{x}_1=5 \cos (100 \pi \mathrm{t}+\pi)(\mathrm{cm})$ và $\mathrm{x}_2=5 \cos (100 \pi \mathrm{t}-\pi / 2)(\mathrm{cm})$. Phương trình dao động tổng hợp của hai dao động trên là
	\choice
	{$x=5 \sqrt{2} \cos \left( 100 \pi t+\dfrac{3\pi}{4}\right)\, (\mathrm{cm})$}
	{\True $x=5 \sqrt{2} \cos \left(100 \pi t-\dfrac{3\pi}{4}\right)\,(\mathrm{cm})$}
	{$x=10 \cos \left(100 \pi t-\dfrac{3\pi}{4}\right)\,(\mathrm{cm})$}
	{$x=10 \cos \left(100 \pi t+\dfrac{3\pi}{4}\right)\,(\mathrm{cm})$}
	\loigiai{
		Ta có 
		\begin{eqnarray*}
			x=x_1+x_2
			&=&5 \cos (100 \pi \mathrm{t}+\pi)+5 \cos (100 \pi \mathrm{t}-\pi / 2)\\
			&=& 5\left[\cos (100 \pi \mathrm{t}+\pi)+\cos (100 \pi \mathrm{t}-\pi / 2)\right]\\
			&=& 5 \cdot 2 cos\left(100 \pi t + \dfrac{\pi}{4} \right) \cos \dfrac{3\pi}{4}   \\
			&=& -5 \sqrt{2} \cos\left(100 \pi t + \dfrac{\pi}{4} \right)\\
			&=& 5 \sqrt{2} \cos \left(100 \pi t-\dfrac{3\pi}{4}\right).
		\end{eqnarray*}
		
	}
\end{ex}\cham{6}

\begin{ex}
	Một vật thực hiện đồng thời hai dao động điều hòa cùng phương, cùng tần số theo các phương trình: $x_1=2 \cos \left(5 \pi t+\dfrac{\pi}{2}\right)(\mathrm{cm})$ ; $x_2=2 \cos (5 \pi t)(\mathrm{cm})$. Biên độ của dao động tổng hợp của hai dao động trên là
	\choice
	{ $2$}
	{ $4$}
	{\True $2\sqrt{2}$}
	{ $\sqrt{2}$}
	\loigiai{
		Xét $x(t)=x_1+x_2=2 \cos \left(5 \pi t+\dfrac{\pi}{2}\right)+2 \cos (5 \pi t)=2\sqrt{2}\cos\left(5\pi t+\dfrac{\pi}{4} \right) $}
\end{ex}\cham{6}
\Closesolutionfile{ans}

\ind{PHẦN II.} \inden{Câu trắc nghiệm đúng sai. Học sinh trả lời từ câu 21 đến câu 25. Trong mỗi ý a), b), c), d) ở mỗi câu, học sinh chọn đúng hoặc sai.}\\
\Opensolutionfile{ans}[ans/1D1-B2-2]

\begin{ex}%[1D1B3-1]
	Cho $\sin\alpha=\dfrac{12}{13}$ và $\dfrac{\pi}{2}<\alpha<\pi$. 
	Xét tính đúng sai của các khẳng định sau:
	\choiceTF
	{\True $\cos \alpha =- \dfrac{5}{13}$}
	{\True $\cos 2\alpha=- \dfrac{119}{169}$}
	{\True $\sin\left(\alpha + \dfrac{\pi}{3}\right)=\dfrac{12 - 5\sqrt{3}}{26}$}
	{\True $\tan\left(\alpha + \dfrac{\pi}{4}\right)= \dfrac{7}{17}$}
	\loigiai{
		\begin{itemchoice}
			\itemch Vì $\dfrac{\pi}{2}<\alpha<\pi$ nên $\cos\alpha= - \sqrt{1 - \sin^2\alpha}= - \sqrt{1 - \left(\dfrac{12}{13}\right)^2}= - \dfrac{5}{13}$.
			\itemch $\cos 2\alpha=2\cos^2\alpha - 1=2\cdot\left( - \dfrac{5}{13}\right)^2 - 1= - \dfrac{119}{169}$.
			\itemch $\sin\left(\alpha + \dfrac{\pi}{3}\right)=\sin\alpha\cos\dfrac{\pi}{3} + \cos\alpha\sin\dfrac{\pi}{3}=\dfrac{12}{13}\cdot\dfrac{1}{2} + \left( - \dfrac{5}{13}\right)\cdot\dfrac{\sqrt{3}}{2}=\dfrac{12 - 5\sqrt{3}}{26}$.
			\itemch Ta có $\tan\alpha=\dfrac{\sin\alpha}{\cos\alpha}=\dfrac{\dfrac{12}{13}}{ - \dfrac{5}{13}}= - \dfrac{12}{5}$.\\
			Suy ra $\tan\left(\alpha + \dfrac{\pi}{4}\right)=\dfrac{\tan\alpha + \tan\dfrac{\pi}{4}}{1 - \tan\alpha\tan\dfrac{\pi}{4}}=\dfrac{ - \dfrac{12}{5} + 1}{1 - \left( - \dfrac{12}{5}\right)\cdot 1}=  \dfrac{7}{17}$.
		\end{itemchoice}
		}
\end{ex}\cham{21}

\begin{ex}%[1T1K3-4]
	Cho $\sin\alpha = \dfrac{3}{5}$, $\cos\beta = \dfrac{12}{13}$ và $0^\circ < \alpha, \beta < 90^\circ$. Xét tính đúng sai của các khẳng định sau:
	\choiceTF
	{$\cos \alpha =-\dfrac{4}{5}$}
	{$\sin \beta =\dfrac{25}{169}$}
	{\True $\sin\left( \alpha + \beta \right) = \dfrac{56}{65}$}
	{\True $\cos\left( \alpha - \beta \right) = \dfrac{63}{65}$}
	\loigiai{
		\begin{itemchoice}
			\itemch Vì $0^\circ < \alpha < 90^\circ$ nên $\cos\alpha > 0$. Do đó
			$$ \cos\alpha = \sqrt{1-\sin^2 \alpha} = \sqrt{1-\left( \dfrac{3}{5} \right)^2} = \dfrac{4}{5}.$$
			\itemch Vì $0^\circ < \beta < 90^\circ$ nên $\sin\beta > 0$. Do đó
			$$ \sin\beta = \sqrt{1-\cos^2 \beta} = \sqrt{1-\left( \dfrac{12}{13} \right)^2} = \dfrac{5}{13}.$$
			\itemch $ \sin\left( \alpha + \beta \right) = \sin\alpha \cos\beta + \cos\alpha \sin\beta = \dfrac{3}{5}\cdot \dfrac{12}{13} + \dfrac{4}{5}\cdot \dfrac{5}{13} = \dfrac{56}{65};$
			\itemch $ \cos\left( \alpha - \beta \right) = \cos\alpha \cos\beta + \sin\alpha \sin\beta = \dfrac{4}{5}\cdot \dfrac{12}{13} + \dfrac{3}{5}\cdot \dfrac{5}{13} = \dfrac{63}{65}.$
		\end{itemchoice}
	}
\end{ex}\cham{16}


\begin{ex}%[1D1B3-1]
	Biết $\cos 2\alpha= - \dfrac{7}{25}$ và $0<\alpha<\dfrac{\pi}{2}$.
	Xét tính đúng sai của các khẳng định sau:
	\choiceTF
	{$\cos \alpha =\dfrac{9}{25}$}
	{$\sin \alpha =\dfrac{16}{25}$}
	{$\tan \alpha =\dfrac{3}{4}$}
	{$\cot \alpha =\dfrac{4}{3}$}
	\loigiai{
		Ta có $\cos 2\alpha= - \dfrac{7}{25}$ nên $2\cos^2\alpha - 1=1 - 2\sin^2\alpha= - \dfrac{7}{25}$
		\begin{itemchoice}
			\itemch Suy ra $\cos^2\alpha=\dfrac{9}{25}$. Với $0<\alpha<\dfrac{\pi}{2}$ thì $\cos \alpha >0$ nên $\cos\alpha=\dfrac{3}{5}$.
			\itemch Suy ra $\sin^2\alpha=\dfrac{16}{25}$. Với $0<\alpha<\dfrac{\pi}{2}$ thì $\sin \alpha >0$ nên $\sin\alpha=\dfrac{4}{5}$.
			\itemch $\tan\alpha=\dfrac{\sin\alpha}{\cos\alpha}=\dfrac{\dfrac{4}{5}}{\dfrac{3}{5}}=\dfrac{4}{3}$.
			\itemch $\cot\alpha=\dfrac{1}{\tan\alpha}=\dfrac{3}{4}$.
		\end{itemchoice}
		}
\end{ex}\cham{16}

\begin{ex}%Câu 30
	\immini{Trên một mảnh giấy hình vuông $ABCD$, bác An đặt một chiếc đèn pin tại vị trí $A$ chiếu chùm sáng phân kì sang góc $C$. Bác An nhận thấy góc chiếu sáng của đèn pin giới hạn bởi hai tia $AM$ và $AN$, ở đó các điểm $M,N$ lần lượt thuộc $BC,CD$ sao cho $BM=\dfrac{1}{2}BC,DN=\dfrac{1}{3}DC$ (Hình bên).
	Xét tính đúng sai của các khẳng định sau:
		\choiceTF
		{\True $\tan \widehat{BAM}=\dfrac{1}{2}$}
		{\True $\tan \widehat{DAN}=\dfrac{1}{3}$}
		{$\widehat{BAM}+\widehat{DAN}=60^\circ $}
		{Góc chiếu sáng của đèn pin bằng $30^\circ$}
	}
	{
		\begin{tikzpicture}
			\def\a{3.5}
			\path 	(0:0) coordinate (A)
			++(0:\a) coordinate (B)
			++(90:\a) coordinate (C)
			($(C)-(B)+(A)$) coordinate (D)
			($(C)!.5!(B)$) coordinate (M)
			($(D)!1/3!(C)$) coordinate (N);
			\draw[thick] 	(A)--(B)--(C)--(D)--cycle
			(M)--(A)--(N);
			\foreach \x/\g in {A/180,B/0,C/45,D/135,M/0,N/90}
			\fill[black] 	(\x) circle (1pt)
			($(\g:3mm)+(\x)$) node {$\x$};	
	\end{tikzpicture}}
	\loigiai{
		 Giả sử hình vuông có độ dài cạnh là $1$.
		 \begin{itemchoice}
		 	\itemch Tam giác $ABM$ vuông tại $B$ có $\tan \widehat{BAM}=\dfrac{BM}{AB}=\dfrac{1}{2}$.
		 	\itemch Tam giác $ADN$ vuông tại $D$ có $\tan \widehat{DAN}=\dfrac{DN}{AD}=\dfrac{1}{3}$.
		 	\itemch $\tan \left( \widehat{BAM}+\widehat{DAN} \right)=\dfrac{\tan \widehat{BAM}+\tan \left( DAN \right)}{1-\tan \widehat{BAM}\tan \left( DAN \right)}=1\Rightarrow \widehat{BAM}+\widehat{DAN}=45^\circ$
		 	\itemch Góc chiếu sáng là $\widehat{MAN}$. Ta có
		 	$$\widehat{MAN}=90^\circ -\left( \widehat{BAM}+\widehat{DAN} \right)=45^\circ $$
		 	Vậy góc chiếu sáng của đèn pin bằng $45^\circ $.
		 \end{itemchoice}
	}
\end{ex}\cham{7}

\begin{ex}%[1D1B3-6]
	Phương trình dao động điều hòa của một vật tại thời điểm $t$ giây được cho bởi công thức $x(t)=A\cos \left(\omega t+\varphi \right)$, trong đó $x(t) $ (cm)  là li độ của vật tại thời điểm $t$ giây, $A$ là biên độ dao động $(A>0)$, $\varphi \in [-\pi; \pi]$ là pha ban đầu của dao động và $\omega$ là tần số góc. Xét hai dao động điều hòa có phương trình lần lượt là
	$$x_1(t)=3\cos\left(\dfrac{\pi}{4}t+\dfrac{\pi}{3}\right)  (\text{cm})\ \text{và} \ x_2(t)=3\cos\left(\dfrac{\pi}{4}t-\dfrac{\pi}{6}\right)  (\text{cm}). $$
	Đặt $x(t)=x_1(t)+x_2(t)$ là phương trình tổng hợp của hai dao động này.
	Xét tính đúng sai của các khẳng định sau:
	\choiceTF
	{\True Tần số góc của dao động tổng hợp là $\dfrac{\pi}{4}$}
	{Phương trình của dao động tổng hợp là $x(t)=3\cos\left(\dfrac{\pi}{4}t+\dfrac{\pi}{12}\right)$}
	{Biên độ của dao động tổng hợp trên bằng $3$}
	{\True Pha ban đầu của dao động tổng hợp trên là $\dfrac{\pi}{12}$ }
	\loigiai{
		\begin{itemchoice}
			\itemch Tần số góc $\omega=\dfrac{\pi}{4}$.
			\itemch Ta có 
			\begin{eqnarray*}
				x(t)&=&x_1(t)+x_2(t)\\
				&=&3\cos\left(\dfrac{\pi}{4}t+\dfrac{\pi}{3}\right) +3\cos\left(\dfrac{\pi}{4}t-\dfrac{\pi}{6}\right)\\
				&=&3.2\cos \dfrac{\left(\dfrac{\pi}{4}t+\dfrac{\pi}{3}\right)+\left(\dfrac{\pi}{4}t-\dfrac{\pi}{6}\right)}{2}\cos \dfrac{\left(\dfrac{\pi}{4}t+\dfrac{\pi}{3}\right)-\left(\dfrac{\pi}{4}t-\dfrac{\pi}{6}\right)}{2}\\
				&=& 6 \cos \dfrac{\dfrac{\pi}{2}t+\dfrac{\pi}{6}}{2} \cos \dfrac{\tfrac{\pi}{2}}{2}\\
				&=&3\sqrt{2}\cos\left(\dfrac{\pi}{4}t+\dfrac{\pi}{12}\right).
			\end{eqnarray*}
			Vậy phương trình của dao động tổng hợp là $x_(t)=3\sqrt{2}\cos\left(\dfrac{\pi}{4}t+\dfrac{\pi}{12}\right).$
			\itemch Căn cứ vào phương trình $x(t)$, ta có biên độ dao động là $3\sqrt{2}$.
			\itemch Căn cứ vào phương trình $x(t)$, ta có pha ban đầu là $\dfrac{\pi}{12}$.
		\end{itemchoice}
	}
\end{ex}\cham{10}

\Closesolutionfile{ans}

\ind{PHẦN III.} \inden{Câu trắc nghiệm trả lời ngắn. Học sinh trả lời từ câu 26 đến câu 30 vào ô kết quả.}\\
\Opensolutionfile{ans}[ans/1D1-B2-3]

\begin{ex}%Câu 15
	Cho hai góc $a$ và $b$ với $\tan a=\dfrac{1}{7}$ và $\tan b=\dfrac{3}{4}$. Khi đó, $\tan \left( a+b \right)$ bằng bao nhiêu?\\
	\shortans[oly]{1}
	\loigiai{
		Ta có $\tan \left( a+b \right)=\dfrac{\tan a+\tan b}{1-\tan a\tan b}$.\\
		Mà $\tan a=\dfrac{1}{7}$ và $\tan b=\dfrac{3}{4}$ nên $\tan \left( a+b \right)=1$.}
\end{ex}\cham{4}

\begin{ex}
	\immini{Cho hình vuông $ABCD$ như hình vẽ bên. Tính $\tan x$ (kết quả làm tròn đến hàng phần trăm).\\
	\shortans[oly]{$0{,}41$}}{
\begin{tikzpicture}[scale=0.75, font=\footnotesize,>=stealth]
	\path
	%	Vẽ mp
	(0,0) coordinate (A)
	(3,0) coordinate (B)
	(3,3) coordinate (C)
	(0,3) coordinate (D)
	(3,1.2) coordinate (E)
	;
	\draw (A)--(B)--(E)node[pos=.5,right]{$5$}--(C)node[pos=.5,right]{$7$}--(D)--(A) (C)--(A)--(E);
	\draw pic["\scriptsize$x$",draw,angle eccentricity=1.6,angle radius=0.5cm]{angle=E--A--C};
	\foreach \x/\g in {A/180,B/0,C/90,D/90,E/0}\draw[fill=black] (\x) circle (.05) +(\g:.5)node{\footnotesize$\x$};
\end{tikzpicture}}
\loigiai{
	Ta có
		\begin{eqnarray*}
		&&\tan\left(45^\circ-x \right)=\tan \widehat{EAB}=\dfrac{5}{12}\\
		&\Leftrightarrow& \dfrac{1-\tan x}{1+\tan x}=\dfrac{5}{12}\\
		&\Leftrightarrow& \tan x = \dfrac{7}{17}\approx 0{,}41.
	\end{eqnarray*}
}
\end{ex}\cham{5}

\begin{ex}
	Biết $\dfrac{\sin \alpha+\sin 2 \alpha}{1+\cos \alpha+\cos 2 \alpha}=n \cdot \tan \alpha$, $n \in \mathbb{N}$. Giá trị $n$ bằng bao nhiêu?\\
	\shortans[oly]{1}
	\loigiai{
		Xét $$\dfrac{\sin \alpha+\sin 2 \alpha}{1+\cos \alpha+\cos 2 \alpha}=\dfrac{\sin \alpha+2\sin  \alpha \cos\alpha}{1+\cos \alpha+2\cos^2 \alpha-1}=\dfrac{\sin \alpha\left(1+2\cos \alpha \right)}{\cos \alpha\left(1+2\cos \alpha \right)}=\dfrac{\sin\alpha}{\sin\alpha}=\tan \alpha.$$
		Suy ra $n=1$.
	}
\end{ex}\cham{5}


\begin{ex}
	\immini{Cho tam giác $ABC$ như hình vẽ. Biết $\widehat{B}=2\widehat{C}=2\alpha$, $AB=2$, $AC=3$. Tính độ dài cạnh $BC$.\\
		\shortans[oly]{$2,5$}}{
		\begin{tikzpicture}[scale=1, font=\footnotesize,>=stealth]
			\path
			%	Vẽ mp
			(1,2) coordinate (A)
			(0,0) coordinate (B)
			(4.4,0) coordinate (C)
			;
			\draw (A)--(B)node[pos=.5,left]{$2$}--(C)node[pos=.5,below]{$x$}--(A)node[pos=.5,right]{$3$};
			\draw pic["\scriptsize$2\alpha$",draw,angle eccentricity=1.6,angle radius=0.5cm]{angle=C--B--A};
			\draw pic["\scriptsize$\alpha$",draw,angle eccentricity=1.6,angle radius=0.5cm]{angle=A--C--B};
			\foreach \x/\g in {A/90,B/180,C/0}\draw[fill=black] (\x) circle (.05) +(\g:.5)node{\footnotesize$\x$};
	\end{tikzpicture}}
\loigiai{

Đặt $BC = x$ ($x > 0$).
Áp dụng định lí sin cho tam giác $ABC$, ta có:
$$ \dfrac{AC}{\sin B} = \dfrac{AB}{\sin C} \Leftrightarrow \dfrac{3}{\sin(2\alpha)} = \dfrac{2}{\sin \alpha} $$
$$ \Leftrightarrow \dfrac{3}{2\sin\alpha\cos\alpha} = \dfrac{2}{\sin\alpha} \Rightarrow \cos\alpha = \dfrac{3}{4}. $$

Áp dụng định lí côsin cho tam giác $ABC$, ta có:
$$ AB^2 = BC^2 + AC^2 - 2 \cdot BC \cdot AC \cdot \cos C $$
$$ \Leftrightarrow 2^2 = x^2 + 3^2 - 2 \cdot x \cdot 3 \cdot \dfrac{3}{4} $$
$$ \Leftrightarrow x^2 - \dfrac{9}{2}x + 5 = 0 \Leftrightarrow \hoac{x &= \dfrac{5}{2} \\ x &= 2.} $$

Kiểm tra nghiệm:
\begin{itemize}
	\item Nếu $x = 2$, tam giác $ABC$ có $AB = BC = 2$ nên cân tại $B$. Khi đó $\widehat{A} = \widehat{C} = \alpha \Rightarrow \widehat{B} = 180^\circ - 2\alpha$.
	Mà giả thiết $\widehat{B} = 2\alpha$ nên $180^\circ - 2\alpha = 2\alpha \Leftrightarrow \alpha = 45^\circ$. Suy ra $\cos \alpha = \dfrac{\sqrt{2}}{2}$ (mâu thuẫn với $\cos\alpha = \dfrac{3}{4}$). Vậy loại $x = 2$.
	\item Nếu $x = \dfrac{5}{2} = 2{,}5$, áp dụng lại định lí côsin ta được $\cos B = \dfrac{BC^2 + AB^2 - AC^2}{2 \cdot BC \cdot AB} = \dfrac{2{,}5^2 + 2^2 - 3^2}{2 \cdot 2{,}5 \cdot 2} = \dfrac{1}{8}$. 
	Lại có $\cos(2C) = 2\cos^2\alpha - 1 = 2\left(\dfrac{3}{4}\right)^2 - 1 = \dfrac{1}{8}$. Suy ra $\cos B = \cos(2C)$, thỏa mãn giả thiết $\widehat{B} = 2\widehat{C}$.
\end{itemize}

Vậy độ dài cạnh $BC$ là $2{,}5$.}
\end{ex}\cham{5}

\begin{ex}
	\immini{Trên nóc một tòa nhà thẳng đứng có một cột ăng-ten cao $5 \mathrm{~m}$. Từ một vị trí quan sát $A$ cao $7 \mathrm{~m}$ so với mặt đất và cách tòa nhà 14 m có thể nhìn thấy đỉnh $B$ và chân $C$ của cột ăng-ten với góc $\widehat{BAC}=10^\circ$. Tính chiều cao của tòa nhà (tính bằng mét và kết quả làm tròn đến hàng đơn vị).\\
		\shortans[oly]{19}
}{
	\begin{tikzpicture}[scale=1, font=\footnotesize,>=stealth]
		\path 
		(0,0) coordinate (A)
		(4,3) coordinate (C)
		(4,4) coordinate (B)
		(4,-1) coordinate (O) 
		(0,-1) coordinate (L) 
		(4,0) coordinate (R) 
		;
		\draw[dashed] (A)--(B) (A)--(C);
		\draw[thick](B)--(C);
		\draw[fill=cyan] (C)--(6,3)--(6,-1)--(4,-1)--(C);
		\foreach \i in {2,3,4,5}
		\draw[fill=white] ({4+1/3},{-1+(2/3)*(\i-1)})--({4+1/3},{-1+(2/3)*(\i-1)+0.5})--({4+2/3},{-1+(2/3)*(\i-1)+0.5})--({4+2/3},{-1+(2/3)*(\i-1)})--cycle 
		({5+1/3},{-1+(2/3)*(\i-1)})--({5+1/3},{-1+(2/3)*(\i-1)+0.5})--({5+2/3},{-1+(2/3)*(\i-1)+0.5})--({5+2/3},{-1+(2/3)*(\i-1)})--cycle;
		\draw[dashed]   (A)--(R)node[pos=.5,below]{$14$ m}; 
		\draw [<->](A)--(L);
		\draw ($(A)!0.5!(L)$) node[left]{$7 \mathrm{~m}$} ($(B)!0.5!(C)$) node[right]{\scriptsize$5 \mathrm{~m}$} ;
		\draw pic[draw=black,double, angle eccentricity=1.4, angle radius=0.6cm]{angle=C--A--B};
		\foreach \x/\g in {A/180,B/140,C/140} 
		\fill[black] (\x) circle (1pt)+(\g:3mm) node {$\x$};
	%	\clip (-0.5,{-1-0.25}) rectangle (6.6,-1) ;
		\draw[fill=gray!30] (-0.5,-1)--(6.6,-1)--(6.6,{-1-0.25})--(-0.5,{-1-0.25})--cycle;
		\foreach \i in {1,2,...,35}
		\draw ({-0.5+0.2*(\i)},-1)--($({-0.5+0.2*(\i)},-1)+(-120:0.3)$) ;
		\node[above] at (0.7,1) {\scriptsize$10^\circ$};
		\draw[->] (0.7,1)--(0.5,0.45);
\end{tikzpicture}}
\loigiai{
	\immini{
	Đặt $CR=h>0$ và $\widehat{CAR}=\alpha$. Ta có
	\begin{itemize}
		\item [$\bullet$] $\tan\alpha=\dfrac{h}{14}$.
		\item [$\bullet$] $\tan\left( \alpha+10^\circ\right) =\dfrac{h+5}{14}$.
	\end{itemize}
	$\tan\left( \alpha+10^\circ\right) =\dfrac{h+5}{14} \Leftrightarrow \dfrac{\tan\alpha + \tan10^\circ}{1-\tan \alpha.\tan10^\circ}=\dfrac{h+5}{14}$.\\
	$\Leftrightarrow \dfrac{\dfrac{h}{14} + \tan10^\circ}{1-\dfrac{h}{14}.\tan10^\circ}=\dfrac{h+5}{14}$. Giải phương trình này, ta tìm được $h \approx 12$. Suy ra, chiều cao ngôi nhà là $h+7=19$ (m).
	}{
	\begin{tikzpicture}[scale=1, font=\footnotesize,>=stealth]
		\path 
		(0,0) coordinate (A)
		(4,3) coordinate (C)
		(4,4) coordinate (B)
		(4,-1) coordinate (O) 
		(0,-1) coordinate (L) 
		(4,0) coordinate (R) 
		;
		\draw[dashed] (A)--(B) (A)--(C);
		\draw[thick](B)--(C);
		\draw[fill=cyan] (C)--(6,3)--(6,-1)--(4,-1)--(C);
		\foreach \i in {2,3,4,5}
		\draw[fill=white] ({4+1/3},{-1+(2/3)*(\i-1)})--({4+1/3},{-1+(2/3)*(\i-1)+0.5})--({4+2/3},{-1+(2/3)*(\i-1)+0.5})--({4+2/3},{-1+(2/3)*(\i-1)})--cycle 
		({5+1/3},{-1+(2/3)*(\i-1)})--({5+1/3},{-1+(2/3)*(\i-1)+0.5})--({5+2/3},{-1+(2/3)*(\i-1)+0.5})--({5+2/3},{-1+(2/3)*(\i-1)})--cycle;
		\draw[dashed]   (A)--(R)node[pos=.5,below]{$14$ m}; 
		\draw [<->](A)--(L);
		\draw ($(A)!0.5!(L)$) node[left]{$7 \mathrm{~m}$} ($(B)!0.5!(C)$) node[right]{\scriptsize$5 \mathrm{~m}$} ;
		\draw pic[draw=black,double, angle eccentricity=1.4, angle radius=0.6cm]{angle=C--A--B};
		\foreach \x/\g in {A/180,B/140,C/140,R/-120} 
		\fill[black] (\x) circle (1pt)+(\g:3mm) node {$\x$};
		%	\clip (-0.5,{-1-0.25}) rectangle (6.6,-1) ;
		\draw[fill=gray!30] (-0.5,-1)--(6.6,-1)--(6.6,{-1-0.25})--(-0.5,{-1-0.25})--cycle;
		\foreach \i in {1,2,...,35}
		\draw ({-0.5+0.2*(\i)},-1)--($({-0.5+0.2*(\i)},-1)+(-120:0.3)$) ;
		\node[above] at (0.7,1) {\scriptsize$10^\circ$};
		\draw[->] (0.7,1)--(0.5,0.45);
\end{tikzpicture}}
	
}
\end{ex}\cham{25}
\Closesolutionfile{ans}


>>>>>>> 33890a9ce1cfbd3079ed36a93e836cf3ef1fde30
